\documentclass[12pt]{article}
\usepackage[utf8]{inputenc}
\usepackage[T1]{fontenc}
\usepackage[margin=1in]{geometry}
\usepackage{setspace}
\usepackage{xcolor}
\usepackage[colorlinks=true, urlcolor=blue!60!black, linkcolor=blue!60!black]{hyperref}
\setlength{\parindent}{0em}
\setlength{\parskip}{0.6em}
\renewcommand{\baselinestretch}{1.0}
\begin{document}
\begin{center}
{\LARGE \textit{Ward v. Rock Against Racism}, 491 U.S. 781 (1989)\par}
\vspace{0.5em}
{\large Case Brief --- Two-Page Edition\par}
\end{center}

\section{Memory Jogger}
A content-neutral time, place, and manner rule must serve a significant interest, not burden substantially more speech than necessary, and leave ample alternatives---but need not be the least restrictive means (\href{https://supreme.justia.com/cases/federal/us/491/781/}{Justia}).

\section{Detailed Case Facts}
Rock Against Racism's Central Park bandshell concerts drew repeated noise complaints, while other sponsors' poor equipment produced inadequate sound. New York City required all bandshell performers to use city equipment run by a city-retained technician, who deferred to sponsors on sound mix.

\section{Procedural History}
The district court upheld the guideline after trial; the Second Circuit reversed, requiring the ``least intrusive'' means.

\section{Judicial Votes}
Reversed, 6--3. Kennedy for the Court (Rehnquist, White, O'Connor, Scalia); Blackmun concurring in the result; Marshall, with Brennan and Stevens, dissenting.

\section{Holding}
Music is protected and the bandshell a public forum, but the guideline is a valid content-neutral regulation. Narrow tailoring is met if the rule promotes a substantial interest that would be served less effectively without it and does not burden substantially more speech than necessary.

\section{Analysis of Opinions}
\textbf{Kennedy:} The government's purpose controls content neutrality, and noise control is unrelated to message. Least-restrictive-means analysis is not part of time, place, and manner review; courts defer to reasonable choices. Only volume was moderated, so alternatives remained. \textbf{Marshall:} the majority ``eviscerat[ed]'' narrow tailoring, and city control of the knob is a prior restraint.

\section{Examples: Future Applications}
\textbf{Same-Side:} A city sound-limiting system at an amphitheater beside a hospital is valid. \textbf{Opposite-Side:} Letting a technician lower volume for ``offensive or hateful lyrics'' is content-based. \textbf{Fence-Sitter:} A total amplification ban on a park's main rally lawn may leave inadequate alternatives.

\section{Critique}
\textit{McCullen v. Coakley} (2014) required proof that less intrusive tools failed, sitting uneasily with \textit{Ward}'s deference. Critics also fault \textit{Ward}'s reliance on unwritten narrowing practices.

\section{Key Quotations}
\begin{quote}
``Lest any confusion on the point remain, we reaffirm today that a regulation of the time, place, or manner of protected speech must be narrowly tailored to serve the government's legitimate, content-neutral interests, but that it need not be the least restrictive or least intrusive means of doing so.''
\end{quote}
\noindent --- Justice Kennedy, for the Court, 491 U.S. at 798--799.

\end{document}
